\section{The weighted divisor graph}\label{sec:graph-setup}

We begin the qualitative argument with the finite-scale graph estimate
that will rule out a correlation bias. Its test functions are arbitrary bounded sequences;
no multiplicativity enters its statement or proof. The arithmetic
application follows in \Cref{sec:analytic-transfer}.

Fix an integer $h\ge1$. Throughout the graph construction use
\begin{equation}\label{eq:graph-constants}
 \eps=10^{-4},\quad \eta=\eps/100,\quad \rho=1/20,\quad
 \kappa=400/\eta,\quad A=\exp(2\kappa),\quad c_*=\eps/100.
\end{equation}
The scale $B$ will be sufficiently large. Define two finite prime sets by
\begin{equation}\label{eq:prime-bands}
 \begin{split}
 \mathcal C&=\{p:B^{1-\eta}<\log p\le B\},\\
 \mathcal Z&=\{p:B^{1-\eps}<\log p\le B^{1-\eta}\},\qquad
 \mathcal S=\mathcal C\cup\mathcal Z,\quad P_0=\exp(B^{1-\eps}).
 \end{split}
\end{equation}
We refer to $\mathcal C$ as the core band and $\mathcal Z$ as the center
band. For $a\in\{C,Z\}$, set $\mathcal P_C=\mathcal C$ and
$\mathcal P_Z=\mathcal Z$, and put
\[
 v_a=\sum_{p\in\mathcal P_a}\frac1p,\qquad
 A_C=A,\quad A_Z=1,\quad \beta_a=(1+A_a)^{-1},\qquad
 \theta=\beta_Z^2=\tfrac14.
\]
Mertens' prime harmonic estimate gives
\[
 v_C=(\eta+o(1))\log B,\qquad
 v_Z=(\eps-\eta+o(1))\log B.
\]

Fix also $\tau\in(1,2)$, $C_0\ge1$, and $T>0$. An admissible divisor
family $\mathcal D$ is any collection of squarefree products of primes
in $\mathcal S$ such that
\begin{equation}\label{eq:divisor-family}
 \begin{gathered}
 d>1,\qquad H<d\le\tau H,\qquad 1\le H\le\exp(C_0B),\\
 \omega(d)\le J:=\lceil C_0\log B\rceil.
 \end{gathered}
\end{equation}
Here $\omega$ counts distinct prime factors. For any integer $x$, define
$\omega_a(x)=\#\{p\in\mathcal P_a:p\mid x\}$, and write
$d_C=\prod_{p\mid d,\ p\in\mathcal C}p$.
For each $d\in\mathcal D$, let $w_d:\Z\to[0,1]$ be any function of
all the residues $x\bmod p$, $p\in\mathcal C$, with support restricted by
\begin{equation}\label{eq:core-cutoff}
 w_d(x)=0\quad\hbox{unless}\quad
 k_d(x):=\#\{p\in\mathcal C:p\mid x,\ p\nmid d\}
       \ge v_C-T\sqrt{v_C}.
\end{equation}
The cutoff need not factor over primes.

For $y=x\pm hd$, define the real edge weight
\begin{equation}\label{eq:graph-kernel}
 g_d(x,y)=w_d(x)w_d(y)A^{\omega_C(d)}\1_{d_C\mid x}
       \prod_{\substack{p\mid d\\p\in\mathcal Z}}
             \left(\1_{p\mid x}-\frac\theta p\right).
\end{equation}
For each $p\mid d$, the residues of $x$ and $y$ agree modulo $p$.
Consequently $g_d(x,y)=g_d(y,x)$. The uncentered core factors enforce
$d_C\mid x$, whereas the center factors can have either sign.

The vertex weight and its mean are
\begin{equation}\label{eq:graph-weights}
 W(x)=\prod_{a\in\{C,Z\}}\beta_a^{-\omega_a(x)},\qquad
 L_0=\prod_{a\in\{C,Z\}}\prod_{p\in\mathcal P_a}(1+A_a/p).
\end{equation}
Let $Q_B=\prod_{p\in\mathcal S}p$. We use the uniform probability space
$\Z/Q_B\Z$, identified by the Chinese remainder theorem with the product
of the uniform spaces $\Z/p\Z$. Its expectation is denoted by $\E$;
$\E_U$ averages only coordinates indexed by $U\subseteq\mathcal S$.
For a uniform residue $n$,
\begin{equation}\label{eq:weight-moments}
 \E W(n)=L_0,\qquad
 \E W(n)^2
 =\prod_a\prod_{p\in\mathcal P_a}
        \left(1+\frac{\beta_a^{-2}-1}{p}\right)
 \le \exp\bigl((A^2+2A)v_C+3v_Z\bigr)=B^{O_A(1)}.
\end{equation}
Both identities follow by independence of the prime coordinates.
The exponent in the last bound is fixed, although large.

\begin{proposition}[Divisor graph estimate]\label{prop:graph-estimate}
Fix $h\ge1$, $1<\tau<2$, $C_0\ge1$, and $T>0$, and use
\eqref{eq:graph-constants}--\eqref{eq:graph-weights}. For every sufficiently
large $B$, every admissible family $\mathcal D$, all permitted cutoffs
$w_d$, all coefficients $|a_d|\le1$, and all functions
$F,G:\N\to\mathbb D$,
\begin{equation}\label{eq:graph-estimate}
 \limsup_{X\to\infty}\frac1X
 \left|\sum_{1\le x\le X}\sum_{d\in\mathcal D}
 a_dF(x)G(x+hd)g_d(x,x+hd)\right|
 \ll L_0 B^{-1-c_*/2}.
\end{equation}
The implied constant and the threshold for $B$ may depend on
$h,\tau,C_0,T$, and the fixed constants in \eqref{eq:graph-constants},
but not on $H,\mathcal D,w_d,a_d,F,G$. All graph parameters are fixed
when $X$ tends to infinity.
\end{proposition}

The proof occupies \Cref{sec:paths,sec:thinning,sec:cylinder-sieve,sec:trace-count,sec:localization}.
The two bands have different roles in that proof. The core weights and
cutoffs give savings from the short divisor interval and from the
lower bound on $k_d(x)$ in \eqref{eq:core-cutoff}.
The center factors provide further cancellation in closed-walk products,
where each edge weight is divided by $W$ at its starting vertex.
The relation $\theta=\beta_Z^2$ matches the
centering constant to that normalization; its precise use is established
in \Cref{sec:trace-count}, once the relevant walk configurations have
been defined.

The comparison scale is $L_0/B$. In the next section, one of $O(B)$
short multiplicative intervals captures enough divisor weight to turn
any persistent correlation bias into a sum of this size with ordinary
divisibility indicators. The extra factor $B^{-c_*/2}$ in
\eqref{eq:graph-estimate} will contradict that
bias once the analytic cost of centering has been bounded.
